% !TeX root = ../main.tex
\subsection{研究目标}

本课题面向三维动态场景，研究连续动态世界模型的构建方法。所建模型由可持续更新的连续场景状态和基于状态的动作条件演化关系组成：利用可用时序视觉观测构建几何、外观及运动状态，依据新增证据调整空间组织和有效结构，再以当前状态及给定动作推演后续变化。研究围绕物体运动、相互遮挡及局部结构变化展开，以代表性具身操作检验所建模型的预测及使用价值。

总体目标是建立由离散观测构建连续状态、组织空间结构、适应局部变化并学习动作条件演化的方法，形成具有明确状态接口的连续动态世界模型研究原型。预期所得状态时间上连续可查询、空间归属稳定，并能在真实局部结构改变后保留有效历史。通过代表性动态场景和选定操作任务，检验各类构建方法对状态恢复及动作后果预测的作用，阐明其成立条件和适用范围。

本课题中的“连续”以场景运动的连续时间表达为基础，使不同采样时刻的观测联系到同一状态过程，并支持状态随新增观测持续更新。结构保持时延续可信联系，发生切割、撕裂等事件时局部调整连接关系和表面集合，保留未受影响区域的有效历史。时空一致性作为状态构建及演化预测的质量要求；研究以可用观测支持的内容为基础，观测不足处保留信息缺失，预测初始状态仅依据当时已经获得的观测和动作历史构建。具体目标如下。

（1）\textbf{构建连续可查询的场景状态。}以离散观测约束场景的连续运动过程，使不同采样时刻的信息在统一的时间表征中建立联系，支持非采样时刻查询和多时刻几何信息累积，为世界模型建立连续时间状态基础。

（2）\textbf{构建空间组织稳定的连续场景状态。}在连续时间状态基础上，以对象或部件组织空间成分，使已有观测和首次显露的表面获得一致的空间归属，支持空间单元内的几何信息累积及运动状态读取，为后续演化提供可组织的空间结构。

（3）\textbf{形成适应拓扑变化的动态表面更新方法。}面向切割、撕裂引起的连接关系改变和物理边界新生，调整失效的结构联系，在延续未变化区域可信信息的基础上更新局部表面，使场景状态能够适应真实结构变化。

（4）\textbf{建立连续状态驱动的动作条件演化模型。}利用连续运动、空间归属及当前有效结构学习给定动作下的场景演化，研究从动态初始状态预测后续变化的方法，检验所得世界模型在预测准确性、一致性及操作中动作选择方面的作用。

\Needspace{110pt}
\subsection{研究内容}

四项研究共同服务于连续动态世界模型的构建。第一项由离散观测构建连续时间状态；第二项利用连续运动建立对象或部件的稳定空间组织；第三项依据观测更新真实结构改变后的有效联系和局部表面；第四项基于所得状态建立动作条件演化关系。前三项形成模型的状态基础，第四项使模型能够推演动作后果。各项分别检验其构建机制，并通过可读取的场景状态逐步衔接。

\subsubsection{连续时间场景状态构建}

动态场景的观测通常按离散时刻采集。以逐时刻状态或相邻帧变换组织运动时，采样时刻之间的演化需要额外的插值或一致性约束，长时间跨度内的几何信息累积也容易受到局部估计误差影响。连续时间建模通过共享的时间函数描述运动过程，使离散观测能够在同一表征中建立联系，为非采样时刻的状态查询和跨时刻几何信息累积提供统一的表达基础。

近年来，动态高斯场景中的连续运动表征逐步发展，已有研究通过时间变形函数、样条轨迹或连续动力学描述场景演化，为连续建模提供了基础\cite{gs4d2024,retimegs2026,particlegs2026}。在复杂动态场景中，连续模型仍面临运动表达能力和观测约束如何匹配的问题：过强的时间约束可能抹去真实的局部变化，过高的运动自由度可能吸收观测噪声，而不同时刻的观测质量和运动变化又会影响几何信息的融合。因此，需要进一步研究连续运动的表达能力同有效观测约束之间的匹配关系。

本项拟研究离散观测约束下的连续时间场景状态构建方法，探索运动表达能力同观测条件相匹配的建模方式，建立多时刻证据对连续运动的有效约束，使几何信息能够沿一致的运动过程累积，形成可查询几何、外观及运动信息的连续状态，作为后续空间组织、结构更新和演化建模的基础。

\subsubsection{连续运动引导的结构化空间表征}

连续时间表征能够描述场景演化，但仍需确定空间成分的归属及信息共享范围。空间邻近或短时运动相似未必意味着同一对象，同一对象内部也可能存在不同部件运动。物体旋转、部件运动或遮挡解除会使已有表面首次进入观测，需要将新观测纳入所属空间单元，完善对已有结构的表征。

已有研究通过运动支撑图、共享运动基和实例信息联系时间运动和空间组织\cite{ref05,shapeofmotion2025,consistentinstance2026}。运动分组同对象身份的关系、观测不充分时的归属可靠性及空间约束对连续状态恢复的作用仍值得研究。本项主要关注物理连接保持时的空间组织，表面首次显露不等同于物理边界新生。

本项拟在连续时间场景状态基础上，研究连续运动证据引导的结构化空间表征方法。结合几何、外观和可见性信息，探索空间归属的建立及延续方式，以对象或部件组织共享运动和局部几何约束，将首次显露的已有表面纳入一致的空间状态，并减少不同单元之间的错误信息传递。预期形成空间归属稳定、能够持续累积几何信息并分别读取对象或部件状态的表示，为世界模型提供有效的空间组织。

\subsubsection{拓扑变化下的动态表面更新}

在连续运动和空间组织基础上，切割、撕裂还会改变物体连接关系，并使内部材料成为新的物体边界。变化前的邻接或共享运动约束可能失效，沿用旧联系容易错误绑定已分离区域，整体放弃历史又会损失稳定区域的有效信息，因此需要适应真实拓扑变化的动态表面表征。

已有方法通过高斯逐帧更新、局部解绑或表面重建适应动态几何变化，为本项提供了直接参考\cite{gaustar2025,atgs2025}。面向真实切割、撕裂，需要进一步研究跨时刻对应、邻接关系及共享运动约束在变化区域的适用范围，并依据物理边界变化更新表面表示。此前未被观测不足以说明表面在物理上新生，区分可见性变化和真实结构变化是避免错误更新的重要条件。

本项拟在结构化场景状态基础上，研究局部变化证据约束的动态表面更新方法。探索历史联系的保留和调整条件，使稳定区域延续可信的运动及几何历史，变化区域解除失效连接，并为真实新生表面建立符合其形成过程的表示。通过在结构事件处局部更新表面集合和连接关系，使连续场景状态能够延续到变化后的有效结构，为后续演化建模提供依据。

\subsubsection{连续状态驱动的动作条件演化建模}

现有面向交互的世界模型围绕图像特征预测、三维场景生成和给定动作下的三维点运动预测等方向开展研究，并将所得预测用于动作规划或策略学习\cite{ref03,gwm2025,ref06}。不同方法分别侧重预训练视觉特征的利用、未来场景的生成或三维交互响应，历史信息和空间先验也已被部分研究引入，为基于场景状态的交互建模提供了基础。

在场景已处于运动、局部观测缺失或物理连接改变的条件下，当前几何相近的状态仍可能具有不同的动作响应。连续运动历史有助于确定动态初始状态，对象和部件归属支持局部交互信息的组织，当前有效连接关系影响局部区域之间的响应关联。已有研究对这些信息的使用方式和支持条件各有侧重，仍需研究如何将持续更新的时空结构状态转化为动作条件预测的有效约束。

本项拟基于前三项形成的连续场景状态，研究动作条件演化关系的构建方法。以截至预测时刻已经恢复的运动和结构为起点，依据对象或部件组织局部交互信息，利用当前有效的连接和边界约束响应预测，学习给定动作下的后续几何和运动变化，使所建世界模型能够从状态恢复进一步推演动作后果。结合新增观测更新状态，并在代表性具身操作任务中检验预测及动作选择的实际效果，明确状态构建的改进能够转化为预测收益的条件。预测初始状态仅依据当时已经获得的观测和动作历史构建。

\Needspace{110pt}
\subsection{拟解决的关键科学问题}

本课题关注的总体问题是：有限离散观测如何约束连续场景状态的可靠构建及其动作条件演化？围绕连续状态构建、稳定空间组织、结构变化适应和动作后果预测，具体提炼以下四个相互关联的科学问题，分别对应四项模型构建内容。

\subsubsection{离散观测对连续场景运动的有效约束}

离散观测在什么条件下能够有效约束连续场景运动？有限采样只能提供部分时刻的状态证据，连续时间形式本身不能保证帧间演化的准确性。需要研究采样条件、观测可靠性及运动复杂程度如何共同影响连续运动恢复，阐明运动表达自由度同有效观测约束之间的匹配关系，以及所得运动描述对多时刻几何信息累积的影响。该问题对应第一项研究内容，重点认识由离散证据构建可靠连续状态的条件。

\subsubsection{连续运动支持稳定空间组织的条件}

连续时间运动证据如何支持稳定的空间组织？局部运动相似和空间邻近不足以单独确定对象或部件的归属，已有表面首次显露时也缺少此前的直接观测对应。需要研究运动过程中的几何、外观和可见性证据如何互补，阐明可靠空间归属的建立条件，以及由此形成的共享约束如何作用于连续运动和几何恢复。该问题对应第二项研究内容，重点认识连续运动支持空间归属形成及延续的条件。

\subsubsection{拓扑变化下结构约束的有效范围}

真实拓扑变化发生时，哪些历史结构联系能够延续，哪些需要调整？表面首次显露和物理边界新生都可能缺少此前的直接观测对应，仅凭对应缺失不足以确定结构已经改变。需要研究局部几何及运动证据如何约束连接失效的判断，阐明结构变化对邻接联系、共享运动约束及表面历史的不同要求，使可信约束在稳定区域延续、在变化区域适应性调整。该问题对应第三项研究内容，重点认识真实结构变化下历史约束的有效范围及其对表面更新的作用。

\subsubsection{连续场景状态对动作后果的有效约束}

连续场景状态在什么条件下能够有效约束动作后果？相近的瞬时几何可能对应不同的运动历史和连接状态，连续可查询或几何精度提高本身不能保证后续响应预测准确。需要研究连续运动、对象或部件组织及有效结构联系如何共同影响动作条件演化，阐明这些信息在观测不充分或局部结构改变时的作用范围，以及状态估计误差如何影响预测可靠性。该问题对应第四项研究内容，重点认识世界模型中连续状态同动作条件演化之间的有效联系及适用条件。
